\section{Concepts of Snort\textregistered}
This section will give a brief introduction to the basic concepts of
Snort\textregistered{}, and discuss these in relation to features made in this
project.

\subsection{Rules}
A rule is basically a set of options that define ways of detecting, classifying
and reporting a potential attack or threat on the network\footnote{A rule
actually consists of a header and options, rtfm.\cite{SnortManual}}.
There are several types of options, but only the general rule options are relevant for this
project. The general rule options provide information about the rule\cite{SnortManual} (i.e.
revision number, alert message). Other types of options, such as payload
options which are actively used in detection, are not relevant for the management of
rules and are thus not considered in this project.

A rule can either be enabled or disabled in the detection system, meaning that
the rule will only have an effect on detection when it is enabled. Since this
project is dealing with multiple sensors connected to one central database, this
functionality has been expanded so that users can keep a rule enabled on one
sensor and disabled on another. In addition, the user can also set an
enabled/disabled flag on rulesets.

A rule \em modifier \em is created for each rule in a ruleset when the ruleset
is added to a sensor. The modifier dictates whether the rule will be enabled or
not on the given sensor, and by default it will assume the original value of
the rule. This means that the original value is only used to determine the
default value of the modifier. A modifier can also be configured to ``follow
the ruleset'', in which case the rule will be enabled only when the ruleset is
enabled, and vice versa. See table \ref{tab:ruleEnabled} for the effects of the
modifier and ruleset options.

%Set row spacing:
\renewcommand{\arraystretch}{1.5}

\begin{table}[tbp]
  \centering
  \begin{tabular}{l|l|l}
  	Modifier & Ruleset & Active on sensor? \\
    \hline
	Enabled & \em any \em & Yes \\
	Disabled & \em any \em & No \\
	Ruleset & Enabled & Yes \\
	Ruleset & Disabled & No \\
  \end{tabular}
  \caption{Determining whether a rule is enabled or disabled on a sensor.}
  \label{tab:ruleEnabled}
\end{table}

\subsection{Generators}
\subsection{Reference Types}
\subsection{Event filters}
Threshold deprecated. Detection filters.